\section{Introduction}
Ce projet vise à construire une bibliothèque de threads en espace utilisateur,
   offrant une interface proche des \texttt{pthreads} tout en s'exécutant sur un
   seul thread noyau. L'intérêt est double : réduire les coûts d'ordonnancement 
  en évitant les transitions noyau, et expérimenter concrètement le changement  
  de contexte via différentes manières (\texttt{ucontext} ou \texttt{setjmp}/\texttt{longjmp}). Ce
  rapport présente nos choix d'implémentation, une comparaison de performances
  entre les différentes variantes, ainsi que les objectifs avancés réalisés.

